\documentclass[10pt,twocolumn]{article}
\usepackage[T1]{fontenc}
\usepackage[utf8]{inputenc}
\usepackage{lmodern}
\usepackage[margin=1.8cm,top=2.4cm,bottom=2cm,columnsep=0.6cm]{geometry}
\usepackage{fancyhdr}
\usepackage{titlesec}
\usepackage{booktabs}
\usepackage{amsmath,amssymb,amsfonts}
\usepackage{microtype}
\usepackage{xcolor}
\usepackage{mdframed}
\usepackage{parskip}
\usepackage{enumitem}
\usepackage{hyperref}

\definecolor{rulered}{RGB}{180,30,30}
\definecolor{boxgray}{RGB}{245,245,245}
\definecolor{keywordgray}{RGB}{100,100,100}

\pagestyle{fancy}
\fancyhf{}
\fancyhead[L]{\textsc{\small Claude Mag}}
\fancyhead[C]{\small\textit{Machine Learning Research Digest}}
\fancyhead[R]{\small 2026-07-22}
\fancyfoot[C]{\small\thepage}
\renewcommand{\headrulewidth}{0.8pt}

\titleformat{\section}{\normalsize\bfseries\color{rulered}}{}{0pt}{}%
  [\vspace{-4pt}\color{rulered}\rule{\columnwidth}{0.4pt}]
\titlespacing{\section}{0pt}{8pt}{4pt}

\hypersetup{colorlinks=true,urlcolor=rulered,linkcolor=black}

\begin{document}
\setlength{\parindent}{0pt}
\setlength{\parskip}{4pt}

{\small\color{keywordgray}\textbf{Keywords:}
  multimodal LLM \textbullet{}
  open weights \textbullet{}
  mixture-of-experts \textbullet{}
  encoder-free \textbullet{}
  thinking mode}\par\vspace{4pt}

{\LARGE\bfseries\raggedright
  Gemma 4 Technical Report}\par\vspace{4pt}

{\large\itshape\raggedright
  Google DeepMind Releases Open-Weight Natively Multimodal Models
  from 2.3B to 31B Parameters, Including an Encoder-Free 12B
  Architecture Ingesting Raw Audio and Image Patches}\par\vspace{6pt}

{\small\textbf{By} Gemma Team (Sherif El Abd et al.)\
  (Google DeepMind)}\par\vspace{2pt}

{\small\color{keywordgray}arXiv:2607.02770 \textbullet{}
  July 2026}
\par\vspace{2pt}\noindent{\color{rulered}\rule{\columnwidth}{0.6pt}}\vspace{6pt}

Google DeepMind's Gemma 4 marks a substantial advance in open-weight
multimodal language models, introducing a family spanning 2.3B to 31B
parameters across dense and Mixture-of-Experts architectures.
The flagship innovation is an encoder-free 12B model that ingests
raw image patches and audio frames through the same token embedding
space as text --- eliminating the modality-specific encoder modules
that have been a fixture of multimodal systems since their inception.
All models include a built-in thinking mode and are released with
full weights, quantized variants, and inference code.

\begin{mdframed}[backgroundcolor=boxgray,linecolor=rulered,linewidth=1pt,
    innerleftmargin=6pt,innerrightmargin=6pt,
    innertopmargin=6pt,innerbottommargin=6pt]
\textbf{\small AT A GLANCE}\par\vspace{4pt}
\begin{description}[leftmargin=0pt,labelsep=4pt,itemsep=2pt,parsep=0pt,
    font=\small\bfseries]
  \item[Problem:] Open-weight multimodal models lag frontier closed
    models on STEM, long-context, and audio-visual tasks; separate
    encoder modules add architectural complexity and maintenance burden.
  \item[Why it matters:] Open-weight multimodal models are essential
    for research reproducibility and on-device deployment; removing
    encoder modules simplifies the architecture fundamentally.
  \item[Method:] Dense (2.3B, 9B), MoE (26B-A4B), and encoder-free
    (12B) architectures; thinking mode trained via SFT + RL; 4-bit QAT
    variants for on-device deployment.
  \item[Result:] Gemma 4 12B (thinking) reaches 72.1\% MMLU-Pro and
    91.4\% MATH-500; preferred over models 2--3$\times$ larger in
    human-rated evaluations.
  \item[Evidence:] Evaluated on MMLU-Pro, MATH-500, MMBench, RULER
    (256K), and speech/audio captioning benchmarks.
\end{description}
\end{mdframed}

\section{The Big Picture}

Open-weight language models have narrowed the gap with frontier closed
models substantially over the past two years, but multimodal open-weight
models have lagged: handling audio natively, processing very long
contexts, and generating reasoning traces have each required
separate architectural components with separate training pipelines.

Gemma 4 unifies these capabilities into a single model family with
a consistent architecture and training procedure.
The encoder-free 12B model is the most architecturally novel member:
by treating image patches and audio frames as tokens in the same
embedding space as text, it eliminates the need for a separately
pretrained and maintained vision encoder and audio encoder,
while retaining competitive performance on both modalities.

\section{Why Existing Approaches Fall Short}

Encoder-based multimodal models couple a visual or audio encoder
(pretrained independently on large corpora) to a language model
through a projection layer.
This approach requires maintaining two (or three) separate model
components with different training schedules, different fine-tuning
protocols, and different quantization sensitivities.
Performance is also bounded by the encoder's representational
quality: if the encoder is undertrained on certain input types
(e.g., audio-visual correspondence), no amount of language model
fine-tuning can recover the missing information.

Thinking modes in prior open-weight models have typically been
implemented via supervised fine-tuning alone, without RL stabilization,
leading to trace inflation --- models generating arbitrarily long
reasoning traces that do not improve answer quality.
Gemma 4 addresses this with a constrained RL stage that optimizes
final-answer correctness under a trace-length penalty.

\section{Core Mechanism}

The encoder-free architecture tokenizes image patches and audio
frames using a learned patchification layer that maps each patch
directly to the model's embedding dimension:
\[
\mathbf{e}_i = W_p \cdot \text{vec}(p_i) + \mathbf{b}_p
\]
where $p_i$ is the $i$-th patch (image or audio), $W_p$ is the
patchification weight matrix, and $\mathbf{b}_p$ is a bias.
The resulting token sequence is concatenated with text tokens and
processed by a standard causal Transformer, enabling end-to-end
attention across all modalities.

\section{Technical Formulation}

The MoE model (26B-A4B) uses top-$k$ expert routing with $k=2$
across $E=64$ experts.
For a given hidden state $\mathbf{h}$, routing scores are computed as:
\[
\mathbf{g}(\mathbf{h}) = \text{softmax}\!\left(\mathbf{h} W_g\right)
\in \mathbb{R}^E
\]
and the top-2 experts are selected and combined:
\[
\text{MoE}(\mathbf{h}) = \sum_{i \in \text{top-2}(\mathbf{g})}
  g_i(\mathbf{h}) \cdot E_i(\mathbf{h})
\]
Load balancing is enforced via an auxiliary loss:
\[
\mathcal{L}_{\text{aux}} = \alpha \sum_{i=1}^{E} f_i \cdot P_i
\]
where $f_i$ is the fraction of tokens routed to expert $i$
and $P_i$ is the mean routing probability to expert $i$.
The thinking mode RL stage applies correctness reward only to
final-answer tokens:
\[
\mathcal{L}_{\text{RL}} = -\mathbb{E}\!\left[R_{\text{correct}} \cdot
  \log \pi_\theta(a_{\text{answer}} \mid c, \tau)\right]
+ \beta \cdot \mathbb{E}[|\tau|]
\]
where $\tau$ is the reasoning trace, $|\tau|$ is trace length,
and $\beta$ is the trace-length penalty.

\section{Learning Procedure}

\begin{enumerate}[itemsep=2pt,parsep=0pt]
  \item \textbf{Pretraining:} All models are pretrained on diverse
    multimodal corpora including web text, code, image-caption pairs,
    and audio-speech pairs.
  \item \textbf{SFT:} Models are fine-tuned on instruction-following
    datasets and on synthetically generated reasoning traces for
    thinking mode initialization.
  \item \textbf{RL:} RLHF with correctness reward for factual and
    mathematical tasks, combined with the trace-length penalty to
    prevent trace inflation.
  \item \textbf{QAT:} Quantization-aware training at 4-bit to
    produce on-device variants with minimal quality loss.
\end{enumerate}

\section{Experimental Findings}

\begin{center}
\small
\begin{tabular}{lcc}
\toprule
Benchmark & Gemma 4 12B & Previous best \\
\midrule
MMLU-Pro & \textbf{72.1\%} & 68.4\% \\
MATH-500 (thinking) & \textbf{91.4\%} & 88.7\% \\
MMBench & \textbf{81.6} & 79.1 \\
RULER 128K & \textbf{97.2\%} & 94.8\% \\
RULER 256K & 87.1\% & --- \\
Speech recog.\ (WER) & \textbf{4.3\%} & 4.5\% \\
\bottomrule
\end{tabular}
\end{center}

The 4-bit QAT variants retain 97--99\% of full-precision benchmark
scores.
Human-rated evaluations prefer Gemma 4 9B and 12B over models with
2--3$\times$ more parameters in head-to-head blind comparisons
on open-ended generation tasks.

\section{Ablations and Interpretation}

\textbf{Encoder-free vs.\ encoder-based:} On standard image
benchmarks, the encoder-free 12B matches the encoder-based 12B
within 1--2 points while reducing total parameters and eliminating
the encoder maintenance burden.

\textbf{Thinking mode cost-benefit:} Enabling thinking mode doubles
average response latency but improves accuracy by 5--15 points on
multi-step reasoning tasks.
It is implemented as an opt-in feature controlled by the user prompt.

\textbf{Trace length penalty:} Without the trace-length penalty in RL,
models generate traces 3--5$\times$ longer with no improvement in
final-answer accuracy --- confirming that unconstrained thinking mode
RL leads to trace inflation rather than deeper reasoning.

\textbf{MoE vs.\ dense at matched active compute:} The 26B-A4B MoE
model outperforms a 4B dense model by 4--6 points across all
benchmarks, confirming that total parameter count provides independent
value beyond active compute.

\section*{Reference}

Gemma Team (Sherif El Abd et al.).\
\textbf{Gemma 4 Technical Report}.
arXiv:2607.02770, July 2026.
\url{https://arxiv.org/abs/2607.02770}

\end{document}
